%% kapa: exam=2 round=24 year=2013 period=3 subject=compensation_law subject_ko=감정평가 및 보상법규 seq=1 number=1 points=40
%% source: Q-Net 감정평가사 기출문제 (cst00309) artlSeq=5077404 file=extracted/1008/5077404/2013년 제24회 감정평가사 2차시험 기출문제/3교시.hwp
%% conversion: kordoc 4.24.51 -> tools/gen2.py (auto); review=proofread-sonnet; flags=none
\begin{longquestion}{【문제 1】}
甲은 S시에 600m\textsuperscript{2}의 토지를 소유하고 있다. S시장 乙은 2002년 5월「국토의 계획 및 이용에 관한 법률」에 의거하여 수립한 도시관리계획으로 甲의 토지가 포함된 일대에 대하여 공원구역으로 지정하였다가 2006년 5월 민원에 따라 甲의 토지를 주거지역으로 변경지정하였다. 乙은 2010년 3월 정부의 녹색도시조성 시책에 부응하여 도시근린공원을 조성하고자 甲의 토지에 대하여 녹지지역으로 재지정하였다. 다음 물음에 답하시오. (40점) \par
1) 甲은 乙이 2010년 3월 그의 토지에 대하여 녹지지역으로 재지정한 것은 신뢰보호의 원칙에 위배될 뿐만 아니라 당해 토지 일대의 이용상황을 고려하지 아니한 결정이었다고 주장하며, 녹지지역 지정을 해제할 것을 요구하고자 한다. 甲의 주장이 법적으로 관철될 수 있는가에 대하여 논하시오. (20점) \par
2) 乙은 공원조성사업을 추진하기 위하여 甲의 토지를 수용하였는데, 보상금산정시 녹지지역을 기준으로 감정평가한 금액을 적용하였다. 그 적법성여부를 논하시오. (20점) \par
\end{longquestion}
